\section{Additional conditions and competence}
Table~\ref{tab:gate} gives the frozen gate checks. Table~\ref{tab:confirmation} applies the competence criteria to the held-out confirmation split. A change in competence between splits leaves the gate-selected population fixed. The original design used separate development, gate, and confirmation splits with 192, 384, and 768 histories, respectively, across both tasks.

\paragraph{Access task.}
Table~\ref{tab:access} gives full-grid estimates for the five confirmed models. In the access task the final clearance maps directly to APPROVE or DENY. Only Gemma 4 12B met confirmation competence in access, passing all three cells. Its gate-selected competing cell (128 histories) gave $D_S=4.384$ [3.480, 5.518] nats and net switching of 0.062 [0.023, 0.109]. Qwen3-8B passed the frozen gate in the interleaved cell (128 histories) but failed confirmation competence there. Its corresponding estimates were 4.221 [3.470, 5.063] nats and 0.273 [0.195, 0.352]. The score intervals are 97.5\%, and switch intervals are 95\%. All other confirmed models failed the access gate in every cell. The confirmation action outputs included one invalid item in the Granite 4.1 8B competing access cell and zero invalid items in routing for all five models.

\paragraph{Balanced histories.}
Table~\ref{tab:conditions} reports all four constructions. An exploratory amendment introduced 192 new routing histories for Qwen3-8B, with identical current states for the target and other entity. The S score shift was 3.895 [3.333, 4.443] nats (95\% interval), and six pairs switched into the wrong-action category and zero out, for a net of 0.031 [0.010, 0.057]. M produced 34 wrong-directed switches and zero reverse switches, for a net of 0.177 [0.125, 0.234]. This was a descriptive follow-up with no prespecified test.

\paragraph{Prediction analysis.}
The secondary H3 analysis asks whether a retrieval-score sensitivity term improves prediction of strict downstream failure. For each entity, measure the paired change in the log-mass margin between the two edited words on the direct current-state query. Subtract the other entity's sensitivity from the target's. Retrieval scoring sums bounded answer-prefix masses across surface forms. The baseline uses pre-decision retrieval scores, word position, prompt length, and design indicators. Both models use fixed L2-regularized logistic regression and five history-grouped folds, with the paired difference in held-out log loss bootstrapped by history. Table~\ref{tab:h3} reports each task separately. Routing gains are inconclusive under this procedure.

\section{Reproducibility}
The confirmation protocol was frozen internally before confirmation. Its revision history includes numerical and answer-event formatting corrections before the final gate. Scoring used unquantized float32 with SDPA on one 80 GB H100, with cached action scores checked against full forwards. Generation-only thinking runs used up to 4,096 new tokens and parsed the answer after the last closing thinking marker. Unclosed, truncated, and unparseable outputs counted as invalid. The balanced follow-up used the confirmation runtime and accepted one small causality deviation under its runner's tolerance. OLMo 3 7B was excluded under a preflight that required exactly zero deviation.

The original analysis used saved score files, manifests, and completion records with hash checks. The tables here use saved derived summaries and their paired intervals. Regenerating the tables from those summaries checks the reporting layer. The sealed input files were unavailable for this build, so raw responses and score provenance could not be independently verified. An anonymous artifact link is omitted.

\begin{table*}[t]
\centering
\begin{tabular}{@{}lcccccc@{}}
\toprule
Role & Main score & Main net & Bal. score & Bal. net & Think 8B & Think 14B \\
\midrule
S & \shortstack{$3.133$\\$[2.794, 3.478]$} & \shortstack{$0.185$\\$[0.143, 0.227]$} & \shortstack{$3.895$\\$[3.333, 4.443]$} & \shortstack{$0.031$\\$[0.010, 0.057]$} & 0/0 & 0/0 \\
O & \shortstack{$1.909$\\$[1.661, 2.161]$} & \shortstack{$0.065$\\$[0.029, 0.099]$} & \shortstack{$2.962$\\$[2.541, 3.365]$} & \shortstack{$0.026$\\$[0.005, 0.052]$} & 0/0 & 2/0 \\
M & \shortstack{$4.807$\\$[4.455, 5.186]$} & \shortstack{$0.294$\\$[0.250, 0.341]$} & \shortstack{$6.621$\\$[5.917, 7.323]$} & \shortstack{$0.177$\\$[0.125, 0.234]$} & 7/0 & 1/0 \\
U & \shortstack{$2.496$\\$[2.201, 2.805]$} & \shortstack{$0.039$\\$[0.013, 0.068]$} & \shortstack{$3.326$\\$[2.889, 3.765]$} & \shortstack{$-0.005$\\$[-0.026, 0.010]$} & 0/0 & 0/0 \\
\bottomrule
\end{tabular}
\caption{Full-grid routing conditions. Main (384 histories) and balanced (192) use Qwen3-8B without thinking. Scores are in nats. The displayed intervals are secondary 95\% intervals. Thinking shows counts into/out of the wrong-action category on 384 main histories per model, with no action scores.}
\label{tab:conditions}
\end{table*}
\begin{table*}[t]
\centering
\begin{tabular}{@{}llcc@{}}
\toprule
Model & Competent cells & $D_S$ (nats) & S net \\
\midrule
Qwen3-8B & none & $3.904$ ${[3.519, 4.306]}$ & $0.167$ ${[0.130, 0.206]}$ \\
Qwen3-14B & none & $4.377$ ${[3.806, 4.960]}$ & $0.161$ ${[0.125, 0.201]}$ \\
Gemma 4 12B & all three & $4.285$ ${[3.699, 4.896]}$ & $0.083$ ${[0.057, 0.112]}$ \\
Qwen3.5 9B & none & $1.444$ ${[1.323, 1.573]}$ & $0.227$ ${[0.185, 0.271]}$ \\
Granite 4.1 8B & none & $3.294$ ${[2.772, 3.840]}$ & $0.068$ ${[0.042, 0.096]}$ \\
\bottomrule
\end{tabular}
\caption{Full-grid access confirmation (384 histories per model, including cells failing competence). Both outcomes use 95\% intervals for these descriptive construction summaries. The gate-selected primary access intervals are reported in the text.}
\label{tab:access}
\end{table*}
\begin{table*}[t]
\centering
\begin{tabular}{@{}lrccccc@{}}
\toprule
Model & Retained & S & O & M & U & M $-$ S \\
\midrule
Qwen3-8B & 335/384 & 0.161 & 0.045 & 0.266 & 0.033 & $0.104$ ${[0.057, 0.155]}$ \\
Qwen3-14B & 359/384 & 0.170 & 0.139 & 0.167 & 0.025 & $-0.003$ ${[-0.050, 0.042]}$ \\
Gemma 4 12B & 339/384 & 0.147 & 0.006 & 0.068 & 0.009 & $-0.080$ ${[-0.124, -0.038]}$ \\
Qwen3.5 9B & 294/384 & 0.228 & 0.085 & 0.177 & 0.044 & $-0.051$ ${[-0.102, 0.000]}$ \\
Granite 4.1 8B & 299/384 & 0.157 & 0.127 & 0.321 & 0.184 & $0.164$ ${[0.117, 0.217]}$ \\
\bottomrule
\end{tabular}
\caption{Net switch rates on the post hoc common tracked routing histories within each model. Paired differences have 95\% intervals.}
\label{tab:tracked}
\end{table*}
\begin{table*}[t]
\centering
\begin{tabular}{@{}lllrrrrrc@{}}
\toprule
Model & Task & Cell & Direct & Current-only & History & Invalid & Easy invalid & Pass \\
\midrule
Qwen3-8B & routing & seq. & 0.977 & 0.938 & 0.779 & 0.000 & 0.000 & yes \\
Qwen3-8B & routing & int. & 0.986 & 0.938 & 0.779 & 0.000 & 0.001 & yes \\
Qwen3-8B & routing & comp. & 0.945 & 0.938 & 0.752 & 0.000 & 0.000 & no \\
Qwen3-8B & access & seq. & 0.889 & 0.906 & 0.822 & 0.000 & 0.055 & no \\
Qwen3-8B & access & int. & 0.953 & 0.938 & 0.877 & 0.000 & 0.013 & yes \\
Qwen3-8B & access & comp. & 0.953 & 0.859 & 0.883 & 0.000 & 0.009 & no \\
Qwen3-14B & routing & seq. & 0.996 & 0.922 & 0.854 & 0.000 & 0.000 & yes \\
Qwen3-14B & routing & int. & 1.000 & 1.000 & 0.846 & 0.000 & 0.000 & yes \\
Qwen3-14B & routing & comp. & 0.997 & 1.000 & 0.814 & 0.000 & 0.000 & yes \\
Qwen3-14B & access & seq. & 0.991 & 0.781 & 0.729 & 0.000 & 0.001 & no \\
Qwen3-14B & access & int. & 0.977 & 0.828 & 0.748 & 0.000 & 0.001 & no \\
Qwen3-14B & access & comp. & 0.946 & 0.828 & 0.697 & 0.000 & 0.000 & no \\
Gemma 3 4B & routing & seq. & 0.939 & 0.688 & 0.645 & 0.000 & 0.023 & no \\
Gemma 3 4B & routing & int. & 0.964 & 0.609 & 0.488 & 0.000 & 0.003 & no \\
Gemma 3 4B & routing & comp. & 0.916 & 0.625 & 0.533 & 0.000 & 0.001 & no \\
Gemma 3 4B & access & seq. & 0.945 & 0.719 & 0.678 & 0.000 & 0.027 & no \\
Gemma 3 4B & access & int. & 0.936 & 0.609 & 0.676 & 0.000 & 0.021 & no \\
Gemma 3 4B & access & comp. & 0.946 & 0.672 & 0.641 & 0.000 & 0.021 & no \\
Ministral 3 14B & routing & seq. & 0.997 & 0.766 & 0.766 & 0.000 & 0.000 & no \\
Ministral 3 14B & routing & int. & 0.986 & 0.766 & 0.756 & 0.000 & 0.000 & no \\
Ministral 3 14B & routing & comp. & 0.999 & 0.812 & 0.756 & 0.000 & 0.000 & no \\
Ministral 3 14B & access & seq. & 0.970 & 0.750 & 0.670 & 0.000 & 0.000 & no \\
Ministral 3 14B & access & int. & 0.926 & 0.812 & 0.711 & 0.000 & 0.003 & no \\
Ministral 3 14B & access & comp. & 0.952 & 0.797 & 0.732 & 0.000 & 0.006 & no \\
Gemma 4 12B & routing & seq. & 1.000 & 0.906 & 0.936 & 0.000 & 0.000 & yes \\
Gemma 4 12B & routing & int. & 0.999 & 0.984 & 0.887 & 0.000 & 0.000 & yes \\
Gemma 4 12B & routing & comp. & 1.000 & 0.969 & 0.930 & 0.000 & 0.000 & yes \\
Gemma 4 12B & access & seq. & 0.912 & 0.969 & 0.967 & 0.000 & 0.088 & no \\
Gemma 4 12B & access & int. & 0.987 & 0.875 & 0.850 & 0.000 & 0.011 & no \\
Gemma 4 12B & access & comp. & 0.993 & 0.906 & 0.932 & 0.000 & 0.003 & yes \\
Qwen3.5 9B & routing & seq. & 1.000 & 0.906 & 0.795 & 0.000 & 0.000 & yes \\
Qwen3.5 9B & routing & int. & 0.922 & 0.922 & 0.746 & 0.000 & 0.000 & no \\
Qwen3.5 9B & routing & comp. & 0.950 & 0.969 & 0.867 & 0.000 & 0.000 & yes \\
Qwen3.5 9B & access & seq. & 0.994 & 0.688 & 0.656 & 0.000 & 0.000 & no \\
Qwen3.5 9B & access & int. & 0.908 & 0.703 & 0.666 & 0.000 & 0.000 & no \\
Qwen3.5 9B & access & comp. & 0.891 & 0.688 & 0.789 & 0.000 & 0.000 & no \\
Granite 4.1 8B & routing & seq. & 0.967 & 0.984 & 0.799 & 0.000 & 0.000 & yes \\
Granite 4.1 8B & routing & int. & 0.955 & 0.953 & 0.738 & 0.000 & 0.000 & yes \\
Granite 4.1 8B & routing & comp. & 0.932 & 0.953 & 0.816 & 0.000 & 0.000 & no \\
Granite 4.1 8B & access & seq. & 0.935 & 0.797 & 0.828 & 0.000 & 0.000 & no \\
Granite 4.1 8B & access & int. & 0.960 & 0.812 & 0.793 & 0.000 & 0.000 & no \\
Granite 4.1 8B & access & comp. & 0.962 & 0.828 & 0.865 & 0.000 & 0.000 & no \\
\bottomrule
\end{tabular}
\caption{Gate competence. Each cell contains 64 histories. Direct is target-current retrieval accuracy. Easy invalid is the fraction of invalid retrieval answers. Both use 11 items per history, as does the action invalid rate. Current-only accuracy uses one action item and historical accuracy uses eight.}
\label{tab:gate}
\end{table*}
\begin{table*}[t]
\centering
\begin{tabular}{@{}lllrrrrrc@{}}
\toprule
Model & Task & Cell & Direct & Current-only & History & Invalid & Easy invalid & Pass \\
\midrule
Qwen3-8B & routing & seq. & 0.988 & 0.953 & 0.858 & 0.000 & 0.000 & yes \\
Qwen3-8B & routing & int. & 0.986 & 0.961 & 0.868 & 0.000 & 0.000 & yes \\
Qwen3-8B & routing & comp. & 0.982 & 0.953 & 0.846 & 0.000 & 0.000 & yes \\
Qwen3-8B & access & seq. & 0.916 & 0.867 & 0.831 & 0.000 & 0.049 & no \\
Qwen3-8B & access & int. & 0.945 & 0.875 & 0.858 & 0.000 & 0.026 & no \\
Qwen3-8B & access & comp. & 0.908 & 0.859 & 0.867 & 0.000 & 0.049 & no \\
Qwen3-14B & routing & seq. & 0.989 & 0.984 & 0.871 & 0.000 & 0.000 & yes \\
Qwen3-14B & routing & int. & 0.996 & 0.977 & 0.851 & 0.000 & 0.000 & yes \\
Qwen3-14B & routing & comp. & 0.997 & 0.992 & 0.843 & 0.000 & 0.000 & yes \\
Qwen3-14B & access & seq. & 0.919 & 0.836 & 0.744 & 0.000 & 0.055 & no \\
Qwen3-14B & access & int. & 0.920 & 0.859 & 0.776 & 0.000 & 0.049 & no \\
Qwen3-14B & access & comp. & 0.962 & 0.820 & 0.747 & 0.000 & 0.024 & no \\
Gemma 4 12B & routing & seq. & 0.999 & 0.914 & 0.898 & 0.000 & 0.001 & yes \\
Gemma 4 12B & routing & int. & 0.997 & 0.898 & 0.840 & 0.000 & 0.001 & no \\
Gemma 4 12B & routing & comp. & 0.977 & 0.914 & 0.905 & 0.000 & 0.022 & yes \\
Gemma 4 12B & access & seq. & 0.969 & 0.969 & 0.937 & 0.000 & 0.030 & yes \\
Gemma 4 12B & access & int. & 0.988 & 0.992 & 0.929 & 0.000 & 0.010 & yes \\
Gemma 4 12B & access & comp. & 0.978 & 0.969 & 0.934 & 0.000 & 0.007 & yes \\
Qwen3.5 9B & routing & seq. & 0.986 & 0.930 & 0.770 & 0.000 & 0.000 & yes \\
Qwen3.5 9B & routing & int. & 0.954 & 0.906 & 0.836 & 0.000 & 0.000 & yes \\
Qwen3.5 9B & routing & comp. & 0.923 & 0.969 & 0.877 & 0.000 & 0.003 & no \\
Qwen3.5 9B & access & seq. & 0.998 & 0.805 & 0.748 & 0.000 & 0.002 & no \\
Qwen3.5 9B & access & int. & 0.997 & 0.789 & 0.759 & 0.000 & 0.000 & no \\
Qwen3.5 9B & access & comp. & 0.965 & 0.805 & 0.846 & 0.000 & 0.009 & no \\
Granite 4.1 8B & routing & seq. & 0.955 & 0.898 & 0.796 & 0.000 & 0.000 & no \\
Granite 4.1 8B & routing & int. & 0.957 & 0.906 & 0.744 & 0.000 & 0.000 & yes \\
Granite 4.1 8B & routing & comp. & 0.963 & 0.930 & 0.790 & 0.000 & 0.000 & yes \\
Granite 4.1 8B & access & seq. & 0.932 & 0.883 & 0.866 & 0.000 & 0.000 & no \\
Granite 4.1 8B & access & int. & 0.920 & 0.828 & 0.824 & 0.000 & 0.000 & no \\
Granite 4.1 8B & access & comp. & 0.949 & 0.867 & 0.893 & 0.001 & 0.000 & no \\
\bottomrule
\end{tabular}
\caption{Confirmation competence. Each cell contains 128 histories. Direct is target-current retrieval accuracy. Easy invalid is the fraction of invalid retrieval answers. Both use 11 items per history, as does the action invalid rate. Current-only accuracy uses one action item and historical accuracy uses eight.}
\label{tab:confirmation}
\end{table*}
\begin{table*}[t]
\centering
\begin{tabular}{@{}llrc@{}}
\toprule
Model & Task & Items & Log-loss gain \\
\midrule
Qwen3-8B & routing & 3072 & $0.0013$ ${[-0.0022, 0.0047]}$ \\
Qwen3-8B & access & 3072 & $-0.0006$ ${[-0.0020, 0.0008]}$ \\
Qwen3-14B & routing & 3072 & $0.0105$ ${[0.0022, 0.0184]}$ \\
Qwen3-14B & access & 3072 & $-0.0001$ ${[-0.0023, 0.0020]}$ \\
Gemma 4 12B & routing & 3072 & $0.0008$ ${[-0.0028, 0.0041]}$ \\
Gemma 4 12B & access & 3072 & $-0.0005$ ${[-0.0024, 0.0012]}$ \\
Qwen3.5 9B & routing & 3072 & $-0.0003$ ${[-0.0025, 0.0019]}$ \\
Qwen3.5 9B & access & 3072 & $0.0038$ ${[-0.0011, 0.0086]}$ \\
Granite 4.1 8B & routing & 3072 & $-0.0000$ ${[-0.0013, 0.0013]}$ \\
Granite 4.1 8B & access & 3072 & $-0.0025$ ${[-0.0050, -0.0004]}$ \\
\bottomrule
\end{tabular}
\caption{H3 on the full confirmation grid (384 histories per model and task). Positive gain means lower out-of-fold log loss with retrieval sensitivity. Paired 95\% history-bootstrap intervals.}
\label{tab:h3}
\end{table*}

